\subsubsection{Stabilité de la "charpente" cristalline : \og critère énergétique \fg{}}

\setcounter{equation}{0}
\setcounter{Q}{0}

Les oxydes \ce{Mn3O4} et \ce{Fe3O4} possèdent une structure spinelle. On va chercher à déterminer l'occupation des sites octaédriques et tétraédriques dans ces deux structures. On sait que la structure spinelle est constituée d'un sous-réseau cubique face centrée d'oxygènes où les cations occupent les sites octaédriques et tétraédriques.

\Q Représenter la maille et en déduire le ratio entre la multiplicité de sites octaédriques et tétraédriques.
\solution{WARNING maille
Le ratio "site O/T" est de 2. On devra prendre en compte ce ratio lors de l'identification des sites occupés par les cations dans un motif spinelle.
}

\Q Donner la configuration électronique de \ce{Mn} et \ce{Fe}.
\solution{D'après les règles de l'Aufbau, on a : 
\begin{align*}
[\ce{Mn}] = ~ _{18}[\ce{Ar}] 3d^5 4s^2 \\
[\ce{Fe}] = ~ _{18}[\ce{Ar}] 3d^6 4s^2
\end{align*}
}

\Q Justifier que ces deux oxydes présentes des \og états de valence mixtes \fg{}. 
\solution{Les oxydes satisfont l'équation de conservation des charges suivantes : 
\begin{align*}
3 no ( \ce{Mn} ) + 4 no ( \ce{O} ) &= 0 \\
3 no ( \ce{Fe} ) + 4 no ( \ce{O} ) &= 0
\end{align*}
On trouve donc les structures suivantes : 
\begin{align*}
\ce{Mn_2^{III} Mn^{II} O4} \\
\ce{Fe_2^{III} Fe^{II} O4}
\end{align*}
}


\Q Rappeler les hypothèses du modèle du champ cristallin et la structure électronique attendue d'un métal du bloc $d$ en configuration O et T.
\solution{Hypothèses : 
\begin{itemize}
\item modèle purement électrostatique ;
\item ligands assimilés à des sphères dures chargées négativement ; 
\item on ne considère que les orbitales $d$ de valence du métal.
\end{itemize}
Il s'agit donc d'un modèle de liaisons ioniques entre le métal et l'oxyde.
}

Dans ces structures, on sait que la configuration électronique imposée par l'anion oxygène est bas champ, haut spin. On négligera l'énergie d'appariement des électrons.

\Q Après avoir défini les termes champs fort/faible et haut/bas spin, prévoir les énergies de stabilisation du champ cristallin (ESCC) en symétrie octaédrique et tétraédrique. Identifier le type de structure normale ou inverse de ces structures.
\solution{Pour certaines configurations de $d^4$ à $d^7$, il existe deux remplissages qui dépendent de la balance entre l'intensité du champ cristallin et le coût énergétique pour apparier les électrons. On a : 
\begin{align*}
\Delta_O > P &\implies \text{champ fort}/\text{bas spin} \\
\Delta_O < P &\implies \text{champ faible}/\text{haut spin}
\end{align*}
La configuration champ faible conduit à un spin du métal plus grand que le champ fort, on qualifie cette configuration de haut spin. On peut calculer l'énergie de stabilisation du champ cristallin à partir des configurations électroniques : 
\begin{center}
\includegraphics[width=0.48\textwidth]{Ex14/Corr01}
\end{center}
Par exemple, calculons l'énergie de champ cristallin du manganèse II dans une configuration octaédrique : 
\begin{align*}
ESCC ( \ce{Mn (II)} )  = 3 \times (-\frac{2}{5} \Delta_O) + 2 \times (\frac{3}{5} \Delta_O) = 0
\end{align*}
On obtient les résultats suivants : 
\begin{center}
\begin{tabular}{ccccc}
\toprule
& \ce{Mn(II)} & \ce{Mn(III)} & \ce{Fe(II)} & \ce{Fe(III)} \\
\midrule
Octaèdre & $0$ & $-3/5 \Delta_O$ & $-2/5 \Delta_O$ & $0$ \\ 
\midrule
\multirow{2}{*}{Tétraèdre} & $0$ & $-2/5 \Delta_T$ & $-3/5 \Delta_T$ & $0$ \\
 & $0$ & $-8/45 \Delta_O$ & $-4/15 \Delta_O$ & $0$ \\
\midrule
$ESCC(O_h) - ESCC (T_d)$ & $0$ & $-19/45 \Delta_O$ & $-2/15 \Delta_O$ & $0$ \\
\bottomrule
\end{tabular}\end{center}
Au bilan, on trouve que Mn(III) et Fe(II) préfèrent une configuration octaédrique. Tandis que les éléments Mn(II) et Fe(III) n'ont aucune préférence entre l'octaèdre et le tétraèdre. Au bilan, on a : 
\begin{align*}
\ce{[Mn_2^{III}]_O [Mn^{II}]_T O4} \quad \ce{[Fe^{III} Fe^{II} ]_O [Fe^{III}]_T  O4}
\end{align*}
}

\Q Expliquer la différence de rayon ionique entre les sites O et T.
\solution{Pour un élément donné, le rayon octaédrique est plus grand que le rayon tétraédrique. Cette différence de taille vient du fait que la longueur de liaison est plus grande dans un octaèdre que dans un tétraèdre.}

\Q En déduire si l'occupation prédit par la théorie du champ cristallin est validée ou invalidée par le critère stérique.
\solution{On peut calculer le rapport $\rho = R_\oplus / R_\ominus$ pour chacun des cations : 
\begin{center}
\begin{tabular}{ccc}
\'{E}lément & $\rho_O$ & $\rho_T$ \\
Mn (II) & 0.69 & 0.57 \\
Mn (III) & 0.56 & -  \\
Fe (II) & 0.66 & 0.55 \\
Fe (III) & 0.56 & 0.45 \\
\end{tabular}
\end{center}
On remarque que l'environnement octaédrique est prédit pour tous les éléments, ce qui est en désaccord avec la théorie du champ cristallin. En réalité, le modèle du champ cristallin prédit la bonne structure pour les deux composés spinelle.}

\paragraph*{Données}
\begin{itemize}
%\item Configurations électroniques : $[Mn] = 3d^5 4s^2$, $[Fe] = 3d^6 4s^2$
\item Numéro atomique : $Z ( \ce{Mn} ) =25$, $Z ( \ce{Fe} ) =26$ ; 
\item Paramètres ioniques des éléments du bloc $d$ : 

\begin{center}
\begin{tabular}{cccc}
\toprule
Cation & Rayon ionique O [\si{\pico\meter}] & Rayon ionique T [\si{\pico\meter}] & $\Delta_O ~ [\si{\per\centi\meter}] $ \\
\midrule
Mn (II) &  97 & 80 & 7500 \\
Mn (III) & 78.5 & - & 21000 \\
Fe (II) & 92 & 77 & 9350 \\
Fe (III) & 78.5 & 63 & 14000 \\
\bottomrule
\end{tabular}
\end{center}

\item On rappelle que $\Delta_T = \frac{4}{9} \Delta_O$;
\item Rayon ionique de l'oxygène au nombre d'oxydation -II : $R_\ominus = 140 ~ \si{\pico\meter}$
\end{itemize}
